% Lösungen Einheit 2 von 4 – Vierecksflächen (zusammenbau v0.1)
\einheitenkopf{Einheit 2 von 4 · Vierecksflächen}

\erg{12}{a) halbes Diagonalenprodukt – es ist eine Raute \quad b) $\overrightarrow{PQ} \cdot \overrightarrow{QR} = 0$, $|PQ| = 5$, $|QR| = 10$, $A = 5 \cdot 10 = 50$ \quad c) $|PR| = 6$, $|QS| = 12$, $A = \tfrac{1}{2} \cdot 6 \cdot 12 = 36$ \quad d) $|PQ| = 8$, $|SR| = 4$, Seitenmitten $M_1(4 | 1 | 1)$ und $M_2(4 | 4 | 5)$, $h = |M_1 M_2| = 5$, $A = \tfrac{1}{2} \cdot (8 + 4) \cdot 5 = 30$ \quad e) $z = 0$ für $t = 2$: $M(6 | 8 | 0)$, $|OM| = 10$; vier rechtwinklige Dreiecke mit den Katheten 10: $A = 4 \cdot \tfrac{1}{2} \cdot 100 = 200$ \quad f) $|PQ| = 5$, $|PS| = k$: $5k = 30$, $k = 6$ \quad g) $P'(0 | 1 + \lambda | 0)$, $S'(0 | 9 - \lambda | 0)$; Trapez mit den parallelen Seiten 8 und $8 - 2\lambda$, Höhe 8: $\tfrac{1}{2} \cdot (16 - 2\lambda) \cdot 8 = 48$, $\lambda = 2$: $P'(0 | 3 | 0)$, $S'(0 | 7 | 0)$}
\erg{13}{a) Die waagerechte Seite wird in wahrer Länge abgebildet. Im Querschnitt senkrecht zu ihr ist die geneigte Seite eine Kathete eines rechtwinkligen Dreiecks, dessen Hypotenuse ihr Schatten auf dem Boden ist; die Hypotenuse ist länger. Also ist das Schattenrechteck größer als das Vordach.}
\erg{14}{a) Mia nimmt den Schenkel $|QR| = 5$ als Höhe. Die Höhe ist der Abstand der parallelen Seiten: $h = 4$, $A = \tfrac{1}{2} \cdot (10 + 4) \cdot 4 = 28$ \quad b) Das Trapez lässt sich in ein Rechteck und Dreiecke zerlegen, deren Flächen alle die Höhe als senkrechten Abstand der Parallelen brauchen; der schräge Schenkel ist als Hypotenuse länger und liefert zu viel. \quad c) $A = 6 \cdot 2 + \tfrac{1}{2} \cdot 6 \cdot 3 = 21$ \quad d) Seitenmitten $(10 | 0 | 0)$ und $(10 | 12 | 9)$, $h = 15\,\mathrm{m}$, $A = \tfrac{1}{2} \cdot (20 + 12) \cdot 15 = 240\,\mathrm{m}^2$, Zuschauer $= 240 : 0{,}6 = 400$}
